% Part: first-order-logic
% Chapter: models-theories
% Section: size-of-structures

\documentclass[../../../include/open-logic-section]{subfiles}

\begin{document}

\olfileid{fol}{mat}{siz}

\olsection{\printtoken{P}{structure} دا سائز ظاہر کرنا}

\begin{explain}
ساختاں دیاں کجھ خاصیتاں اسیں زبان دیاں غیر منطقی علامتاں
ورتے بغیر وی ظاہر کر سکدے آں۔ مثال لئی کجھ
!!{sentence}s ایداں دے نیں جیہڑے کسے !!a{structure}
وچ جے تے صرف جے سچے ہوندے نیں جدوں اس !!{structure}
دے !!{domain} وچ !!{element}s دی گنتی کسے خاص
مثبت عدد~$n$ توں گھٹ نہ ہووے، ودھ نہ ہووے، یا عین اوہنی ہووے۔
% PNB-SEM-026: نقطیاں والے خاکے مثبت گنتی مطابق محدود جوڑاں دے مخفف نیں؛ اک لئی خالی جوڑ سچ اے۔
ہیٹھاں نقطیاں والیاں عبارتیں مطلوب گنتی مطابق محدود جوڑاں
تے متبادلاں دے مخفف نیں۔ اک عنصر لئی جوڑی وار وکھرے ہون
دیاں شرطاں دی خالی جوڑ-فہرست نوں سچ سمجھیا جاندا اے۔
\end{explain}

\begin{prop}
!!{sentence}
\begin{multline*}
  !A_{\ge n} \ident \lexists[x_1][\lexists[x_2][\dots\lexists[x_n][{}]]]\\
  \begin{aligned}
  (\eq/[x_1][x_2] \land {}
  \eq/[x_1][x_3] \land \eq/[x_1][x_4] \land \dots \land \eq/[x_1][x_n] \land {}\\
\eq/[x_2][x_3] \land \eq/[x_2][x_4] \land \dots \land {} \eq/[x_2][x_n] \land {} \\
\vdots\\
\eq/[x_{n-1}][x_n])
\end{aligned}
\end{multline*}
!!a{structure}~$\Struct M$ وچ جے تے صرف جے سچا اے
جدوں $\Domain M$ وچ گھٹ توں گھٹ $n$ !!{element}s
ہون۔ ایس لئی $\Sat{M}{\lnot !A_{\ge n+1}}$ جے تے
صرف جے جدوں $\Domain M$ وچ ودھ توں ودھ~$n$
!!{element}s ہون۔

\end{prop}

\begin{prop}
!!{sentence}
\begin{multline*}
!A_{= n} \ident \lexists[x_1][\lexists[x_2][\dots\lexists[x_n][{}]]] \\
\begin{aligned}
  (\eq/[x_1][x_2] \land {}
  \eq/[x_1][x_3] \land \eq/[x_1][x_4] \land \dots \land \eq/[x_1][x_n] \land {}\\
\eq/[x_2][x_3] \land \eq/[x_2][x_4] \land \dots \land {} \eq/[x_2][x_n] \land {} \\
\vdots\\
\eq/[x_{n-1}][x_n] \land {} \\
\lforall[y][(\eq[y][x_1] \lor \dots \lor \eq[y][x_n]]))
\end{aligned}
\end{multline*}
!!a{structure}~$\Struct M$ وچ جے تے صرف جے سچا اے
جدوں $\Domain M$ وچ عین $n$ !!{element}s ہون۔
\end{prop}

\begin{prop}
اک !!{structure} لامتناہی اے جے تے صرف جے اوہ
\[
\{!A_{\ge 1}, !A_{\ge 2}, !A_{\ge 3}, \dots \}.
\]
دا نمونہ اے۔
\end{prop}

ایداں دا کوئی اک خالص منطقی جملہ نہیں جیہڑا~$\Struct
M$ وچ جے تے صرف جے سچا ہووے جدوں $\Domain M$
لامتناہی ہووے۔ البتہ غیر منطقی !!{predicate}s ورتن
والے ایداں دے !!{sentence}s دتے جا سکدے نیں جنہاں
دے صرف لامتناہی نمونے ہون (بھاویں ہر لامتناہی
!!{structure} انہاں دا نمونہ نہ ہووے)۔ محدود ساخت
ہون دی خاصیت !!{sentence}s دے لامتناہی سیٹ نال وی ظاہر
نہیں کیتی جا سکدی۔
% PNB-SEM-027: ناقابلِ شمار ہون بارے ایہ دعویٰ ایتھے ورتی گننے جوگی زبان لئی اے۔
ایتھے ورتی گننے جوگی زبان وچ !!{nonenumerable} ساخت ہون دی
خاصیت وی جملیاں دے لامتناہی سیٹ نال ظاہر
نہیں کیتی جا سکدی۔ ایہ حقیقتاں محدود گواہی دے
قضیے تے لووینہائم--سکولم (L\"owenheim--Skolem)
دے قضیاں توں نکلدیاں نیں۔

\end{document}
